% Recuperación de 02_trit_geometrias y de la holonomía con memoria.
\subsubsection{Parametrización temporal de la compatibilidad y orientación trítica}
\label{trit-temporal-compatibilidad}

La interpretación temporal del TRIT se establece sobre las transiciones del
estado completo. Sea \(x\xrightarrow{\rm adm}y\) la relación de admisibilidad determinada
por las operaciones del TPK en \(\widehat X\). Una trayectoria parametrizada
por un conjunto discretamente ordenado \(\mathcal T\) es una aplicación
\[
 \gamma:\mathcal T\longrightarrow\widehat X,\qquad
 \gamma(t)\xrightarrow{\rm adm}\gamma(t^+),
\]
donde \(t^+\) es el sucesor de \(t\). La relación de compatibilidad precede a
la elección del parámetro: el orden permite expresar sus transiciones como
una evolución. Cuando se utiliza el término sección, se declara además la
proyección \(p:E\to\mathcal T\) y la condición
\(p\circ\gamma=\operatorname{id}_{\mathcal T}\). Así quedan especificados
el espacio de estados, el orden temporal y la información que se transporta.

La firma \(\tau\) selecciona un régimen cuadrático y admite una lectura
temporal relativa a esa parametrización. El tipo \(+1\), con
\(J_{+1}^{2}=-I\), representa retorno y memoria; el tipo \(0\), con
\(J_0^{2}=0\), representa el umbral de transición; el tipo \(-1\), con
\(J_{-1}^{2}=I\), representa apertura en dos direcciones propias. En la
terminología temporal del modelo, estas funciones corresponden,
respectivamente, a pasado, presente y futuro. La correspondencia se refiere
a la posición operacional de una transición dentro de una historia:
antecedente conservado, actualización y prolongación admisible.

Esta interpretación adquiere contenido mediante tres operaciones distintas.
La restricción de una trayectoria recupera el prefijo ya construido. La
actualización \(\mathcal U_t\) compone selección, transporte y registro en
la transición observada. La prolongación considera las continuaciones que
conservan las condiciones de frontera de ese prefijo. Si
\(\mathcal P_n(x)\) denota el conjunto de continuaciones admisibles de
longitud \(n\) desde \(x\), la supresión del último paso define
\[
 r_n:\mathcal P_{n+1}(x)\longrightarrow\mathcal P_n(x).
\]
Una continuación ilimitada es una colección \((\gamma_n)_{n\geq1}\) que
satisface \(r_n(\gamma_{n+1})=\gamma_n\). El sistema inverso conserva de
este modo las restricciones que cada profundidad impone sobre las
anteriores. La existencia de una colección compatible se establece en el
dominio de prolongación correspondiente; el mero establecimiento de los
conjuntos \(\mathcal P_n(x)\) no sustituye esa demostración.

La bidireccionalidad combina, por tanto, la reconstrucción de antecedentes
con la compatibilidad de continuaciones. Sobre la sucesión de firmas, la
involución \(C_{\rm rev}\) invierte el orden y cambia el signo de cada
componente. Sobre la trayectoria actúan, además, las transformaciones de
posición, hoja y frontera especificadas por el transporte. La firma
invertida es una observación de la trayectoria conjugada, no un reemplazo
de sus restantes coordenadas. El umbral permanece distinguido: el valor
visible \(0\) conserva su función de transición incluso cuando su cociente
o su registro histórico son distintos de cero.

La relación con las geometrías exponenciales es igualmente precisa. En
cada régimen fijo, \(\exp(tJ_\tau)\) compone parámetros y conserva la
norma algebraica. El cierre elíptico pertenece a esa representación local.
El retorno de fase de \(\Gamma_9\), en cambio, actúa sobre la trayectoria
con memoria, conforme a \eqref{eq:holonomia-memoria}. Una trayectoria puede
recuperar su fase y conservar un nuevo antecedente: el dato histórico ha
cambiado aunque la observación de fase sea la misma. El umbral parabólico
retiene un desplazamiento lineal; la apertura hiperbólica separa dos
direcciones propias de expansión y contracción. Estas tres realizaciones
proporcionan leyes de composición, mientras que el TPK determina su
selección y su concatenación efectiva.

La suspensión aperiódica de la sección \ref{sec:generacion} hará explícita
esta distinción mediante un factor de fase finito, una rotación irracional
y un grado de memoria. El vínculo entre TRIT y temporalidad queda así
expresado por operaciones sobre trayectorias y por sus representaciones
cuadráticas. Su contenido excede una clasificación de signos: conserva el
régimen, el sentido del transporte y la relación entre antecedente,
transición y continuación.
